\chapter{Geração de Código Intermediário}

\section{Introdução}

A cadeia de transformações de um compilador raramente salta diretamente da
árvore sintática abstrata (AST) para o código de máquina.  Tal salto seria
possível, mas produziria código de baixa qualidade e tornaria o compilador
fortemente acoplado à arquitetura-alvo.  A estratégia mais usada é
introduzir uma
\textbf{representação intermediária} (RI), uma linguagem interna de grau de
abstração situado entre a linguagem-fonte de alto nível e o código de máquina.

A figura a seguir resume o posicionamento da RI no \emph{pipeline}:

\begin{center}
  \begin{tikzpicture}[
      node distance=1.2cm and 1.4cm,
      box/.style={draw, rounded corners, minimum width=2.6cm, minimum
        height=0.9cm,
      font=\small\sffamily, align=center},
      arr/.style={-Stealth, thick}
    ]
    \node[box] (ast)  {AST anotada};
    \node[box, right=of ast]  (ir)   {RI};
    \node[box, right=of ir]   (opt)  {RI otimizada};
    \node[box, right=of opt]  (asm)  {Código alvo};
    \draw[arr] (ast) -- (ir)  node[midway, above,
    font=\tiny\sffamily]{tradução};
    \draw[arr] (ir)  -- (opt) node[midway, above,
    font=\tiny\sffamily]{otimizações};
    \draw[arr] (opt) -- (asm) node[midway, above, font=\tiny\sffamily]{seleção};
  \end{tikzpicture}
\end{center}

Este capítulo define formalmente a RI utilizada na disciplina, apresenta sua
semântica operacional \emph{big-step} e descreve as regras de
\textbf{tradução dirigida por sintaxe} que convertem a AST das linguagens
TWhile e TImp para essa RI.

\section{Propriedades Desejáveis}\label{sec:ir-desiderata}

Uma boa RI deve satisfazer as seguintes propriedades:

\begin{description}
  \item[Simplicidade.] Poucos construtores facilitam a análise e a
    transformação do código.  Um IR com apenas $\approx 15$ formas sintáticas
    é preferível a um que espelhe todas as construções da linguagem-fonte.
  \item[Independência de máquina.] A RI não deve codificar detalhes do
    hardware-alvo (convenção de chamada, número de registradores, tamanho de
    ponteiro).  Isso permite que um \emph{front-end} gere RI para
    múltiplas arquiteturas.
  \item[Independência de linguagem.] A RI não deve pressupor construções
    específicas da linguagem-fonte, de modo que diferentes \emph{front-ends}
    possam compartilhar o mesmo \emph{back-end}.  LLVM IR e JVM bytecode são
    exemplos reais desse princípio.
  \item[Suporte a transformações.] A RI deve facilitar análises de fluxo e
    otimizações como propagação de constantes, eliminação de subexpressões
    comuns e movimentação de código invariante a laços.
  \item[Proximidade do código-alvo.] A RI deve expor operações de memória e
    controle de fluxo explícitos, de modo que a seleção de instruções seja
    direta.
\end{description}

\section{Abordagens Principais}\label{sec:ir-abordagens}

Existem quatro famílias de RI amplamente usadas:

\begin{enumerate}
  \item \textbf{RI em árvore de baixo nível.} A AST é rebaixada
    (\emph{lowered}) para uma árvore com operações explícitas de memória e
    saltos incondicionais/condicionais.  É a abordagem adotada neste capítulo
    e no compilador \textbf{Tiger} de Appel.
  \item \textbf{Código de máquina de pilha (\emph{bytecode}).} Instruções
    operam em uma pilha de operandos virtual.  Exemplos: \emph{p-code}
    (Pascal), JVM, CPython.
  \item \textbf{Código de três endereços / quádruplas.} Cada instrução tem a
    forma $t \leftarrow a \oplus b$, onde $t$, $a$ e $b$ são temporários ou
    constantes.  Organizado em blocos básicos e grafo de fluxo de controle.
    Exemplo: LLVM IR.
  \item \textbf{Representação em estilo de passagem de continuações (CPS).}
    Comum em compiladores de linguagens funcionais.  Todo salto torna-se uma
    chamada de função; não há implicitamente um ``próximo passo''.
\end{enumerate}

Neste capítulo adotamos a abordagem (1): uma \textbf{representação
intermediária em árvore} (doravante IRT).

\section{A Representação Intermediária em Árvore (IRT)}\label{sec:irt}

A IRT distingue dois tipos de nós: \textbf{expressões} (que produzem um valor
de tamanho de palavra) e \textbf{comandos} (que produzem efeitos colaterais).

\subsection{Sintaxe das Expressões}\label{sec:irt-expr}

\begin{Definition}[Expressões da IRT]\label{def:irt-expr}
  O conjunto das expressões da IRT é dado pela gramática:
  \begin{align*}
    e \;::=\; & \mathbf{CONST}(n)                       & n \in \mathbb{Z} \\
    \mid\; & \mathbf{TEMP}(t)                    & t \in \mathit{Temp} \\
    \mid\; & \mathbf{NAME}(\ell)                 & \ell \in \mathit{Label} \\
    \mid\; & \mathbf{BINOP}(\mathit{op}, e_1, e_2) \\
    \mid\; & \mathbf{MEM}(e)                     \\
    \mid\; & \mathbf{CALL}(e_f, e_1, \ldots, e_n) \\
    \mid\; & \mathbf{ESEQ}(s, e)
  \end{align*}
  onde $\mathit{op} \in \{
    \mathtt{ADD}, \mathtt{SUB}, \mathtt{MUL}, \mathtt{DIV}, \mathtt{MOD},
    \mathtt{AND}, \mathtt{OR},  \mathtt{XOR},
    \mathtt{LSHIFT}, \mathtt{RSHIFT}, \mathtt{ARSHIFT},
    \mathtt{EQ},  \mathtt{NEQ},
    \mathtt{LT},  \mathtt{LEQ}, \mathtt{GT},  \mathtt{GEQ}
  \}$.
\end{Definition}

A semântica informal de cada construtor é:

\begin{center}
  \renewcommand{\arraystretch}{1.3}
  \begin{tabular}{>{\ttfamily}l p{9.5cm}}
    \toprule
    \normalfont Construtor & Significado \\
    \midrule
    CONST(n)            & Constante inteira $n$.  Booleanos são codificados como
    $0$ (falso) e $1$ (verdadeiro). \\
    TEMP(t)             & Valor armazenado no temporário $t$.  Temporários são
    variáveis sem endereço; o alocador de registradores os
    mapeará para registradores ou posições na pilha. \\
    NAME($\ell$)        & Endereço de memória associado ao rótulo
    $\ell$.  Permite
    passar o endereço de uma função ou de um rótulo de salto
    como dado. \\
    BINOP(op, $e_1$, $e_2$) & Aplica a operação binária \textit{op}
    aos valores de
    $e_1$ e $e_2$.  Operações de comparação produzem $1$
    (verdadeiro) ou $0$ (falso). \\
    MEM($e$)            & Lê a palavra de memória no endereço $e$.  No lado
    esquerdo de \textbf{MOVE}, denota escrita. \\
    CALL($e_f$, $\overline{e}$) & Chama a função cujo endereço é $e_f$, com
    argumentos $\overline{e}$.  O resultado fica em
    registradores de retorno convencionais. \\
    ESEQ($s$, $e$)      & Executa o comando $s$ pelos seus efeitos e
    então avalia
    $e$ como valor da expressão inteira.  Permite embutir
    efeitos dentro de uma expressão. \\
    \bottomrule
  \end{tabular}
\end{center}

\subsection{Sintaxe dos Comandos}\label{sec:irt-stmt}

\begin{Definition}[Comandos da IRT]\label{def:irt-stmt}
  O conjunto dos comandos da IRT é dado por:
  \begin{align*}
    s \;::=\; & \mathbf{MOVE}(e_{\mathit{dst}}, e)
    & e_{\mathit{dst}} \in \{\mathbf{TEMP}(t),\; \mathbf{MEM}(e')\} \\
    \mid\; & \mathbf{EXP}(e) \\
    \mid\; & \mathbf{SEQ}(s_1, s_2) \\
    \mid\; & \mathbf{JUMP}(e) \\
    \mid\; & \mathbf{CJUMP}(e, \ell_t, \ell_f) \\
    \mid\; & \mathbf{LABEL}(\ell) \\
    \mid\; & \mathbf{RETURN}(e_1, \ldots, e_n)
  \end{align*}
\end{Definition}

\begin{center}
  \renewcommand{\arraystretch}{1.3}
  \begin{tabular}{>{\ttfamily}l p{9.5cm}}
    \toprule
    \normalfont Construtor & Significado \\
    \midrule
    MOVE(TEMP($t$), $e$) & Atribui o valor de $e$ ao temporário $t$. \\
    MOVE(MEM($e_a$), $e$) & Escreve o valor de $e$ no endereço de
    memória $e_a$. \\
    EXP($e$)             & Avalia $e$ e descarta o resultado; usado para efeitos
    colaterais (e.g., chamadas sem valor de retorno). \\
    SEQ($s_1$, $s_2$)    & Executa $s_1$ e depois $s_2$ sequencialmente. \\
    JUMP($e$)            & Salta incondicionalmente para o endereço $e$. \\
    CJUMP($e$, $\ell_t$, $\ell_f$) & Se $e \neq 0$ salta para $\ell_t$;
    caso contrário, para $\ell_f$. \\
    LABEL($\ell$)        & Define o rótulo $\ell$ neste ponto do código.
    Não produz instrução de máquina; apenas marca uma
    posição endereçável. \\
    RETURN($\overline{e}$) & Encerra a função corrente devolvendo os valores
    $\overline{e}$ ao chamador. \\
    \bottomrule
  \end{tabular}
\end{center}

\begin{remark}
  A restrição sobre o destino de \textbf{MOVE} (apenas \textbf{TEMP} ou
  \textbf{MEM}) é intencional: ela evita destinos arbitrários e facilita a
  geração de código, pois toda escrita em memória ou em registrador é
  sintaticamente distinguível.
\end{remark}

\section{Semântica Operacional Big-Step da IRT}\label{sec:irt-semantics}

\subsection{Estado e Valores}\label{sec:irt-state}

\begin{Definition}[Estado]\label{def:irt-state}
  Um \textbf{estado} é um par $\sigma = \langle \tau, \mu \rangle$ onde:
  \begin{itemize}
    \item $\tau : \mathit{Temp} \rightharpoonup \mathbb{Z}$ é o
      \emph{armazenamento de temporários}, uma função parcial de temporários
      em inteiros;
    \item $\mu : \mathbb{Z} \rightharpoonup \mathbb{Z}$ é a
      \emph{memória}, que mapeia endereços inteiros em palavras.
  \end{itemize}
  Escrevemos $\sigma.\tau$ e $\sigma.\mu$ para acessar os componentes.
  As notações $\tau[t \mapsto v]$ e $\mu[a \mapsto v]$ denotam,
  respectivamente, a atualização do temporário $t$ e da posição de memória $a$
  com o valor $v$.
\end{Definition}

Todos os valores da IRT são inteiros de precisão de máquina ($\mathbb{Z}$).
Ponteiros, booleanos e outros tipos são codificados como inteiros.

Um \textbf{programa IRT} é uma função parcial
$P : \mathit{Label} \rightharpoonup \mathit{Stmt}$
que mapeia cada rótulo a um comando (ou sequência de comandos).

\subsection{Semântica das Expressões}\label{sec:irt-expr-sem}

A relação $\sigma \vdash e \Downarrow v$ significa ``no estado $\sigma$, a
expressão $e$ avalia para o inteiro $v$''.

\subsubsection*{Constantes e Temporários}

\[
  \infer[\textsc{Const}]
  {\sigma \vdash \mathbf{CONST}(n) \Downarrow n}
  {}
  \qquad
  \infer[\textsc{Temp}]
  {\sigma \vdash \mathbf{TEMP}(t) \Downarrow v}
  {\sigma.\tau(t) = v}
\]

\[
  \infer[\textsc{Name}]
  {\sigma \vdash \mathbf{NAME}(\ell) \Downarrow \mathit{addr}(\ell)}
  {}
\]

A função $\mathit{addr}(\ell)$ devolve o endereço de código associado ao
rótulo $\ell$; assume-se que este mapeamento é fixado em tempo de link.

\subsubsection*{Operações Binárias}

\[
  \infer[\textsc{BinOp}]
  {\sigma \vdash \mathbf{BINOP}(\mathit{op}, e_1, e_2) \Downarrow
  \mathit{op}(v_1, v_2)}
  {\sigma \vdash e_1 \Downarrow v_1 \quad
  \sigma \vdash e_2 \Downarrow v_2}
\]

A interpretação de $\mathit{op}(v_1, v_2)$ é a habitual.
Para operadores de comparação ($\mathtt{EQ}$, $\mathtt{LT}$, etc.),
o resultado é $1$ se a relação se mantém e $0$ caso contrário:
\[
  \mathtt{EQ}(v_1, v_2) =
  \begin{cases} 1 & \text{se } v_1 = v_2 \\ 0 & \text{caso contrário.}
  \end{cases}
\]

\subsubsection*{Acesso à Memória}

\[
  \infer[\textsc{Mem}]
  {\sigma \vdash \mathbf{MEM}(e) \Downarrow v}
  {\sigma \vdash e \Downarrow a \quad \sigma.\mu(a) = v}
\]

\subsubsection*{Chamada de Função}

\[
  \infer[\textsc{Call}]
  {\sigma \vdash \mathbf{CALL}(e_f, e_1, \ldots, e_n) \Downarrow v}
  {\sigma \vdash e_f   \Downarrow a \quad
    \sigma \vdash e_i   \Downarrow v_i \quad
  \mathit{invoke}(a, v_1,\ldots,v_n) = v}
\]

A função auxiliar $\mathit{invoke}$ modela a invocação da função residente no
endereço $a$ com os argumentos fornecidos.  Seu resultado é o valor de
retorno.  A implementação concreta de $\mathit{invoke}$ depende da convenção
de chamada e está fora do escopo desta apresentação formal.

\subsubsection*{Expressão Sequenciada}

\[
  \infer[\textsc{ESeq}]
  {\sigma \vdash \mathbf{ESEQ}(s, e) \Downarrow v}
  {P, \sigma \vdash s \Downarrow \langle \sigma', \blacksquare \rangle \quad
  \sigma' \vdash e \Downarrow v}
\]

O $\mathbf{ESEQ}$ executa o comando $s$ (modificando o estado) e depois
avalia $e$ no estado resultante.  A notação $P, \sigma \vdash s \Downarrow
\langle \sigma', \blacksquare \rangle$ é definida na seção seguinte.

\subsection{Semântica dos Comandos}\label{sec:irt-stmt-sem}

A relação $P, \sigma \vdash s \Downarrow r$ significa
``dado o programa $P$ e o estado inicial $\sigma$, o comando $s$ executa e
produz o resultado $r$''.

\begin{Definition}[Resultado de Execução]\label{def:irt-result}
  O resultado $r$ de um comando é:
  \[
    r \;::=\; \langle \sigma', \blacksquare \rangle
    \;\mid\; \langle \sigma', \uparrow (v_1,\ldots,v_k) \rangle
  \]
  O símbolo $\blacksquare$ denota terminação normal e
  $\uparrow(v_1,\ldots,v_k)$ denota retorno de função com valores
  $v_1,\ldots,v_k$.
\end{Definition}

\subsubsection*{Atribuição a Temporário}

\[
  \infer[\textsc{MoveTemp}]
  {P, \sigma \vdash \mathbf{MOVE}(\mathbf{TEMP}(t), e)
  \Downarrow \langle \sigma', \blacksquare \rangle}
  {\sigma \vdash e \Downarrow v \quad
  \sigma' = \langle \sigma.\tau[t \mapsto v],\; \sigma.\mu \rangle}
\]

\subsubsection*{Escrita em Memória}

\[
  \infer[\textsc{MoveMem}]
  {P, \sigma \vdash \mathbf{MOVE}(\mathbf{MEM}(e_a), e)
  \Downarrow \langle \sigma', \blacksquare \rangle}
  {\sigma \vdash e_a \Downarrow a \quad
    \sigma \vdash e   \Downarrow v \quad
  \sigma' = \langle \sigma.\tau,\; \sigma.\mu[a \mapsto v] \rangle}
\]

\subsubsection*{Expressão como Comando}

\[
  \infer[\textsc{Exp}]
  {P, \sigma \vdash \mathbf{EXP}(e) \Downarrow \langle \sigma,
  \blacksquare \rangle}
  {\sigma \vdash e \Downarrow \_}
\]

O valor produzido por $e$ é descartado; apenas os efeitos colaterais (como
uma chamada de função com efeitos) são relevantes.

\subsubsection*{Sequência}

\[
  \infer[\textsc{SeqNorm}]
  {P, \sigma \vdash \mathbf{SEQ}(s_1, s_2)
  \Downarrow \langle \sigma'', r \rangle}
  {P, \sigma  \vdash s_1 \Downarrow \langle \sigma', \blacksquare \rangle \quad
  P, \sigma' \vdash s_2 \Downarrow \langle \sigma'', r \rangle}
\]

\[
  \infer[\textsc{SeqRet}]
  {P, \sigma \vdash \mathbf{SEQ}(s_1, s_2)
  \Downarrow \langle \sigma', \uparrow \overline{v} \rangle}
  {P, \sigma \vdash s_1 \Downarrow \langle \sigma', \uparrow
  \overline{v} \rangle}
\]

Se $s_1$ retorna, $s_2$ não é executado: o resultado de retorno propaga-se
diretamente.

\subsubsection*{Rótulo}

\[
  \infer[\textsc{Label}]
  {P, \sigma \vdash \mathbf{LABEL}(\ell) \Downarrow \langle \sigma,
  \blacksquare \rangle}
  {}
\]

$\mathbf{LABEL}$ não produz efeito em tempo de execução; é uma anotação
estática para o vinculador e para os saltos.

\subsubsection*{Salto Incondicional}

\[
  \infer[\textsc{Jump}]
  {P, \sigma \vdash \mathbf{JUMP}(\mathbf{NAME}(\ell))
  \Downarrow r}
  {P(\ell) = s \quad P, \sigma \vdash s \Downarrow r}
\]

O salto consulta $P$ para obter a continuação do rótulo $\ell$ e a executa.

\subsubsection*{Salto Condicional}

\[
  \infer[\textsc{CJumpTrue}]
  {P, \sigma \vdash \mathbf{CJUMP}(e, \ell_t, \ell_f) \Downarrow r}
  {\sigma \vdash e \Downarrow v \quad v \neq 0 \quad
  P(\ell_t) = s \quad P, \sigma \vdash s \Downarrow r}
\]

\[
  \infer[\textsc{CJumpFalse}]
  {P, \sigma \vdash \mathbf{CJUMP}(e, \ell_t, \ell_f) \Downarrow r}
  {\sigma \vdash e \Downarrow 0 \quad
  P(\ell_f) = s \quad P, \sigma \vdash s \Downarrow r}
\]

\subsubsection*{Retorno}

\[
  \infer[\textsc{Return}]
  {P, \sigma \vdash \mathbf{RETURN}(e_1, \ldots, e_k)
  \Downarrow \langle \sigma, \uparrow (v_1, \ldots, v_k) \rangle}
  {\sigma \vdash e_i \Downarrow v_i \quad (1 \leq i \leq k)}
\]

\begin{remark}
  A semântica big-step apresentada não termina para programas que executam
  laços infinitos, o que é esperado: a semântica big-step é definida apenas
  para computações que terminam.  Para raciocinar sobre não-terminação,
  utiliza-se semântica small-step ou a semântica de rastreamentos
  (\emph{trace semantics}).
\end{remark}

\section{Tradução Dirigida por Sintaxe}\label{sec:irt-sdt}

A tradução de TWhile e TImp para a IRT é definida por duas funções de
tradução indutivas sobre a AST:

\begin{itemize}
  \item $\mathcal{E}\llbracket e \rrbracket$ : traduz uma expressão
    de entrada em
    uma expressão IRT;
  \item $\mathcal{S}\llbracket s \rrbracket$ : traduz um comando de
    entrada em um
    comando IRT.
\end{itemize}

Assumimos que a AST foi previamente verificada pelo analisador semântico, de
modo que todos os tipos são conhecidos e corretos.

\subsection{Tradução de Expressões}\label{sec:sdt-expr}

\subsubsection*{Literais}

\begin{align*}
  \mathcal{E}\llbracket \mathtt{false} \rrbracket &= \mathbf{CONST}(0) \\
  \mathcal{E}\llbracket \mathtt{true}  \rrbracket &= \mathbf{CONST}(1) \\
  \mathcal{E}\llbracket n \rrbracket              &= \mathbf{CONST}(n)
  \qquad n \in \mathbb{Z}
\end{align*}

Strings são tratadas como ponteiros para dados estáticos e são aqui omitidas
por brevidade.

\subsubsection*{Variáveis}

\[
  \mathcal{E}\llbracket x \rrbracket = \mathbf{TEMP}(x)
\]

Cada variável local é mapeada para um temporário de mesmo nome.  O alocador
de registradores decide posteriormente se o temporário vive em um registrador
ou em uma posição na pilha.

\subsubsection*{Operações Aritméticas e Relacionais}

\[
  \mathcal{E}\llbracket e_1 \oplus e_2 \rrbracket =
  \mathbf{BINOP}(\widehat{\oplus},\;
    \mathcal{E}\llbracket e_1 \rrbracket,\;
  \mathcal{E}\llbracket e_2 \rrbracket)
\]

onde $\widehat{\oplus}$ é o operador IRT correspondente:

\begin{center}
  \renewcommand{\arraystretch}{1.2}
  \begin{tabular}{cc|cc|cc}
    \toprule
    TWhile & IRT & TWhile & IRT & TWhile & IRT \\
    \midrule
    \texttt{+}  & ADD & \texttt{==} & EQ  & \texttt{<}  & LT  \\
    \texttt{-}  & SUB & \texttt{!=} & NEQ & \texttt{<=} & LEQ \\
    \texttt{*}  & MUL & \texttt{>}  & GT  & \texttt{>=} & GEQ \\
    \texttt{/}  & DIV & \texttt{\&\&} & AND & \texttt{||} & OR \\
    \texttt{\%} & MOD & & & & \\
    \bottomrule
  \end{tabular}
\end{center}

\subsubsection*{Negação Unária e \texttt{not}}

\begin{align*}
  \mathcal{E}\llbracket -\, e \rrbracket
  &= \mathbf{BINOP}(\mathtt{SUB},\; \mathbf{CONST}(0),\;
  \mathcal{E}\llbracket e \rrbracket) \\
  \mathcal{E}\llbracket \mathtt{not}\; e \rrbracket
  &= \mathbf{BINOP}(\mathtt{EQ},\;
  \mathcal{E}\llbracket e \rrbracket,\; \mathbf{CONST}(0))
\end{align*}

\subsubsection*{Chamada de Função}

\[
  \mathcal{E}\llbracket f(e_1, \ldots, e_n) \rrbracket =
  \mathbf{CALL}\!\left(\mathbf{NAME}(f),\;
    \mathcal{E}\llbracket e_1 \rrbracket,\; \ldots,\;
  \mathcal{E}\llbracket e_n \rrbracket\right)
\]

\subsection{Tradução de Comandos}\label{sec:sdt-stmt}

\subsubsection*{Declaração de Variável Local}

\[
  \mathcal{S}\llbracket \mathtt{let}\; x : T := e \mathtt{;} \rrbracket =
  \mathbf{MOVE}(\mathbf{TEMP}(x),\; \mathcal{E}\llbracket e \rrbracket)
\]

\subsubsection*{Atribuição Simples}

\[
  \mathcal{S}\llbracket x := e \mathtt{;} \rrbracket =
  \mathbf{MOVE}(\mathbf{TEMP}(x),\; \mathcal{E}\llbracket e \rrbracket)
\]

\subsubsection*{Retorno}

\[
  \mathcal{S}\llbracket \mathtt{return}\; e \mathtt{;} \rrbracket =
  \mathbf{RETURN}(\mathcal{E}\llbracket e \rrbracket)
\]

\[
  \mathcal{S}\llbracket \mathtt{return} \mathtt{;} \rrbracket =
  \mathbf{RETURN}()
\]

\subsubsection*{Sequência de Comandos}

\[
  \begin{aligned}
    \mathcal{S}\llbracket s_1\; s_2\; \cdots\; s_n \rrbracket
    &= \mathbf{SEQ}(\mathcal{S}\llbracket s_1 \rrbracket, \\
      &\quad \mathbf{SEQ}(\mathcal{S}\llbracket s_2 \rrbracket, \\
    &\qquad \cdots\; \mathcal{S}\llbracket s_n \rrbracket \cdots))
  \end{aligned}
\]

\subsection{Tradução de Estruturas de Controle}\label{sec:sdt-control}

\subsubsection*{Condicional}

Para traduzir $\mathtt{if}\; e\; \{\, b_{\mathit{then}} \,\}$,
geramos dois novos rótulos $\ell_t$ e $\ell_f$:

\[
  \begin{aligned}
    \mathcal{S}\llbracket \mathtt{if}\; e\; \{b_t\} \rrbracket
    &= \mathbf{SEQ}( \\
      &\quad \mathbf{CJUMP}(\mathcal{E}\llbracket e \rrbracket,\;
      \ell_t,\; \ell_f), \\
      &\quad \mathbf{SEQ}( \\
        &\qquad \mathbf{LABEL}(\ell_t), \\
        &\qquad \mathbf{SEQ}( \\
          &\qquad\quad \mathcal{S}\llbracket b_t \rrbracket, \\
    &\qquad\quad \mathbf{LABEL}(\ell_f))))
  \end{aligned}
\]

Para $\mathtt{if}\; e\; \{b_t\}\; \mathtt{else}\; \{b_f\}$, introduzimos um
rótulo adicional $\ell_{\mathit{end}}$:

\[
  \begin{aligned}
    \mathcal{S}\llbracket \mathtt{if}\; e\; \{b_t\}\; \mathtt{else}\;
    \{b_f\} \rrbracket
    &= \mathbf{SEQ}( \\
      &\quad \mathbf{CJUMP}(\mathcal{E}\llbracket e \rrbracket,\;
      \ell_t,\; \ell_f), \\
      &\quad \mathbf{SEQ}( \\
        &\qquad \mathbf{LABEL}(\ell_t), \\
        &\qquad \mathbf{SEQ}( \\
          &\qquad\quad \mathcal{S}\llbracket b_t \rrbracket, \\
          &\qquad\quad \mathbf{SEQ}( \\
            &\qquad\qquad \mathbf{JUMP}(\mathbf{NAME}(\ell_{\mathit{end}})), \\
            &\qquad\qquad \mathbf{SEQ}( \\
              &\qquad\qquad\quad \mathbf{LABEL}(\ell_f), \\
              &\qquad\qquad\quad \mathbf{SEQ}( \\
                &\qquad\qquad\qquad \mathcal{S}\llbracket b_f \rrbracket, \\
    &\qquad\qquad\qquad \mathbf{LABEL}(\ell_{\mathit{end}})))))))
  \end{aligned}
\]

O $\mathbf{JUMP}$ ao final do ramo \emph{then} evita que a execução
``caia'' no corpo do ramo \emph{else}.

\subsubsection*{Laço \texttt{while}}

Para $\mathtt{while}\; e\; \{b\}$, introduzimos três novos rótulos:
$\ell_h$ (cabeça do laço), $\ell_b$ (corpo) e $\ell_f$ (após o laço):

\[
  \begin{aligned}
    \mathcal{S}\llbracket \mathtt{while}\; e\; \{b\} \rrbracket
    &= \mathbf{SEQ}( \\
      &\quad \mathbf{LABEL}(\ell_h), \\
      &\quad \mathbf{SEQ}( \\
        &\qquad \mathbf{CJUMP}(\mathcal{E}\llbracket e \rrbracket,\;
        \ell_b,\; \ell_f), \\
        &\qquad \mathbf{SEQ}( \\
          &\qquad\quad \mathbf{LABEL}(\ell_b), \\
          &\qquad\quad \mathbf{SEQ}( \\
            &\qquad\qquad \mathcal{S}\llbracket b \rrbracket, \\
            &\qquad\qquad \mathbf{SEQ}( \\
              &\qquad\qquad\quad \mathbf{JUMP}(\mathbf{NAME}(\ell_h)), \\
    &\qquad\qquad\quad \mathbf{LABEL}(\ell_f))))))
  \end{aligned}
\]

O $\mathbf{JUMP}(\mathbf{NAME}(\ell_h))$ ao final do corpo fecha o laço,
retornando ao teste da condição.

\subsection{Tradução de Funções}\label{sec:sdt-fn}

Uma declaração de função de nível superior

\[
  \mathtt{fn}\; f\;
  (p_1 : T_1,\; \ldots,\; p_n : T_n) : T_r\; \{b\}
\]

é traduzida em um \textbf{fragmento de função} $(f,\; s_{\mathit{body}})$
onde:

\[
  \begin{aligned}
    s_{\mathit{body}}
    &= \mathbf{SEQ}( \\
      &\quad \mathbf{MOVE}(\mathbf{TEMP}(p_1),\; \mathbf{TEMP}(a_1)), \\
      &\quad \cdots \\
      &\quad \mathbf{SEQ}( \\
        &\qquad \mathbf{MOVE}(\mathbf{TEMP}(p_n),\; \mathbf{TEMP}(a_n)), \\
    &\qquad \mathcal{S}\llbracket b \rrbracket))
  \end{aligned}
\]

Os temporários $a_1, \ldots, a_n$ representam os registradores (ou
posições de pilha) onde a convenção de chamada deposita os argumentos.
O prólogo da função copia cada argumento para o temporário do parâmetro
formal correspondente.

\section{Exemplo Completo}\label{sec:irt-example}

Consideremos a função fatorial:

\begin{minted}{text}
fn fat(n : int) : int {
    if n <= 1 {
        return 1;
    } else {
        return n * fat(n - 1);
    }
}
\end{minted}

Aplicando as regras de tradução, obtemos o fragmento IRT (omitindo os
rótulos numéricos por brevidade):

\begin{minted}{text}
fat:
  MOVE(TEMP(n), TEMP(a0))           -- copia arg da conv. de chamada
  CJUMP(BINOP(LEQ, TEMP(n), CONST(1)), L_then, L_else)
  LABEL(L_then)
    RETURN(CONST(1))
  LABEL(L_else)
    RETURN(BINOP(MUL,
             TEMP(n),
             CALL(NAME(fat),
                  BINOP(SUB, TEMP(n), CONST(1)))))
\end{minted}

Observe que cada ramo do \texttt{if-else} termina com \textbf{RETURN}, de
modo que não é necessário um rótulo de confluência ao final.

\section{Formas Canônicas}\label{sec:irt-canonical}

A IRT recém-gerada contém dois tipos de ``ruído'' que complicam as fases
posteriores:

\begin{enumerate}
  \item \textbf{$\mathbf{ESEQ}$ em posições arbitrárias.}
    $\mathbf{ESEQ}$ pode reordenar efeitos colaterais de maneira
    não-óbvia quando aparece dentro de operandos de $\mathbf{BINOP}$ ou
    $\mathbf{CALL}$.
  \item \textbf{$\mathbf{CALL}$ dentro de expressões.}
    O valor de retorno de uma chamada precisa ser capturado em um temporário
    antes de ser usado como operando.
\end{enumerate}

O processo de \textbf{canonicalização} elimina esses problemas em duas
passagens:

\begin{description}
  \item[Linearização.] $\mathbf{ESEQ}$ é ``içado'' (\emph{hoisted}) para
    o nível de comando via reescrita:
    \[
      \begin{aligned}
        &\mathbf{BINOP}(\mathit{op},\; \mathbf{ESEQ}(s, e_1),\; e_2)
        \;\longrightarrow\;
        \mathbf{ESEQ}( \\
          &\quad s, \\
        &\quad \mathbf{BINOP}(\mathit{op},\; e_1,\; e_2))
      \end{aligned}
    \]
    desde que $s$ não contenha efeitos que interfiram com $e_2$.

  \item[Extração de chamadas.] Toda $\mathbf{CALL}$ dentro de uma expressão
    é substituída por
    \[
      \begin{aligned}
        &\mathbf{ESEQ}( \\
          &\quad \mathbf{MOVE}(\mathbf{TEMP}(t),\; \mathbf{CALL}(\cdots)), \\
        &\quad \mathbf{TEMP}(t))
      \end{aligned}
    \]
    para um novo temporário $t$.
\end{description}

Após a canonicalização, a IRT está em \textbf{forma linear}: uma sequência
plana de comandos sem $\mathbf{ESEQ}$, que pode ser facilmente particionada
em \textbf{blocos básicos} para análise de fluxo de controle.

\section{Implementação em Haskell: Geração de Código para
TWhile}\label{sec:ir-impl-haskell}

Esta seção descreve a implementação concreta do \emph{backend} de geração de
código intermediário para a linguagem TWhile, desenvolvida em Haskell no módulo
\texttt{TWhile.Backend.IR.TWhileCodegen}.  A implementação segue de perto as
regras de tradução apresentadas na Seção~\ref{sec:irt-sdt} e serve como
referência direta entre a teoria e o código executável.

\subsection{Mônada de Geração de Código}\label{sec:ir-impl-monad}

A geração de código requer dois efeitos distintos:

\begin{enumerate}
  \item \textbf{Geração de rótulos:} os construtores \texttt{SWhile}
    e \texttt{SIf} necessitam de rótulos únicos para os pontos de entrada
    e saída de cada bloco de controle.  Um contador global garante a unicidade.
  \item \textbf{Sinalização de erros:} situações inesperadas durante a
    compilação são propagadas como mensagens de erro sem abortar o programa.
\end{enumerate}

Esses dois efeitos são combinados via transformadores de mônada:

\begin{minted}{haskell}
data CgState = CgState
  { nextLabel :: Int
  , stmtAcc   :: [IR.Stmt]
  }

type CgM a = StateT CgState (ExceptT String Identity) a
\end{minted}

O estado \texttt{CgState} mantém:
\begin{itemize}
  \item \texttt{nextLabel}: o próximo inteiro livre para formar um rótulo único;
  \item \texttt{stmtAcc}: a lista de instruções IRT emitidas até o momento
    (o acumulador de saída).
\end{itemize}

A camada \texttt{ExceptT String Identity} provê propagação de erros sem I/O;
\texttt{Identity} é usada porque a geração de código é uma transformação pura.
Dois auxiliares encapsulam as operações mais comuns:

\begin{minted}{haskell}
freshLabel :: String -> CgM IR.Label
freshLabel prefix = do
  n <- gets nextLabel
  modify (\s -> s { nextLabel = n + 1 })
  return (prefix ++ "_" ++ show n)

emit :: IR.Stmt -> CgM ()
emit s = modify (\st -> st { stmtAcc = stmtAcc st ++ [s] })
\end{minted}

\texttt{freshLabel} consome o contador e devolve um rótulo como
\texttt{"while\_head\_0"}, \texttt{"if\_true\_1"}, etc.
\texttt{emit} acrescenta uma instrução ao acumulador.

\subsection{Ponto de Entrada}\label{sec:ir-impl-entry}

A função principal recebe a AST de um programa TWhile e produz um
\texttt{IR.Program}, isto é, uma lista de \texttt{FuncDef}:

\begin{minted}{haskell}
compileTWhile :: TWhile -> Either String IR.Program
compileTWhile (TWhile stmts) =
  case runCgM (mapM_ codegenStmt stmts) of
    Left err -> Left err
    Right ss -> Right [IR.FuncDef "main" [] (foldStmts ss)]
\end{minted}

O programa TWhile inteiro é embutido em uma única função \texttt{main} sem
parâmetros.  A função auxiliar \texttt{foldStmts} converte a lista plana de
instruções em uma árvore \texttt{SEQ} associada à direita:

\begin{minted}{haskell}
foldStmts :: [IR.Stmt] -> IR.Stmt
foldStmts []     = IR.EXP (IR.CONST 0)
foldStmts [s]    = s
foldStmts (s:ss) = IR.SEQ s (foldStmts ss)
\end{minted}

\subsection{Geração de Comandos}\label{sec:ir-impl-stmts}

A função \texttt{codegenStmt} implementa as regras de tradução de comandos.
Cada caso consome um construtor da AST de TWhile e emite zero ou mais
instruções IRT.

\subsubsection*{Declaração e Atribuição}

\begin{minted}{haskell}
codegenStmt (SDecl v _ e) = do
  ir <- codegenExp e
  emit (IR.MOVE (IR.TEMP v) ir)

codegenStmt (SAssign v e) = do
  ir <- codegenExp e
  emit (IR.MOVE (IR.TEMP v) ir)
\end{minted}

Ambos os casos traduzem diretamente para
$\mathbf{MOVE}(\mathbf{TEMP}(v),\; \mathcal{E}\llbracket e \rrbracket)$,
conforme as regras da Seção~\ref{sec:sdt-stmt}.  A anotação de tipo na
declaração é descartada, pois a IRT é não-tipada.

\subsubsection*{Impressão e Leitura}

\begin{minted}{haskell}
codegenStmt (SPrint e) = do
  ir <- codegenExp e
  emit (IR.EXP (IR.CALL (IR.NAME "print") [ir]))

codegenStmt (SRead _ v) =
  emit (IR.MOVE (IR.TEMP v) (IR.CALL (IR.NAME "read_int") []))
\end{minted}

A impressão compila para uma chamada à função embutida \texttt{print}, cujo
efeito colateral é exibir o valor no terminal.
A leitura descarta a expressão de \emph{prompt} (que seria uma string, não
representável na IRT inteira) e obtém o valor via a função embutida
\texttt{read\_int}.

\subsubsection*{Laço \texttt{while}}

\begin{minted}{haskell}
codegenStmt (SWhile e body) = do
  lh <- freshLabel "while_head"
  lb <- freshLabel "while_body"
  lf <- freshLabel "while_exit"
  emit (IR.LABEL lh)
  cond <- codegenExp e
  emit (IR.CJUMP cond lb lf)
  emit (IR.LABEL lb)
  mapM_ codegenStmt body
  emit (IR.JUMP (IR.NAME lh))
  emit (IR.LABEL lf)
\end{minted}

A implementação segue fielmente o esquema da Seção~\ref{sec:sdt-control}:

\begin{center}
  \renewcommand{\arraystretch}{1.2}
  \begin{tabular}{ll}
    \toprule
    Instrução emitida & Papel \\
    \midrule
    \texttt{LABEL lh}            & Cabeça do laço; destino do salto
    de retorno \\
    \texttt{CJUMP(cond, lb, lf)} & Avalia a condição; entra no corpo ou sai \\
    \texttt{LABEL lb}            & Início do corpo \\
    \textit{corpo}               & Instruções geradas recursivamente \\
    \texttt{JUMP(NAME lh)}       & Fecha o laço, volta ao teste \\
    \texttt{LABEL lf}            & Ponto de saída do laço \\
    \bottomrule
  \end{tabular}
\end{center}

\subsubsection*{Condicional \texttt{if-else}}

\begin{minted}{haskell}
codegenStmt (SIf e b1 b2) = do
  lt <- freshLabel "if_true"
  lf <- freshLabel "if_false"
  le <- freshLabel "if_end"
  cond <- codegenExp e
  emit (IR.CJUMP cond lt lf)
  emit (IR.LABEL lt)
  mapM_ codegenStmt b1
  emit (IR.JUMP (IR.NAME le))
  emit (IR.LABEL lf)
  mapM_ codegenStmt b2
  emit (IR.LABEL le)
\end{minted}

O \texttt{JUMP(NAME le)} ao final do ramo \emph{then} é essencial: sem ele a
execução ``cairia'' para o ramo \emph{else} após concluir o ramo verdadeiro.

\subsection{Geração de Expressões}\label{sec:ir-impl-exprs}

A função \texttt{codegenExp} implementa $\mathcal{E}\llbracket \cdot
\rrbracket$ e devolve um \texttt{IR.Expr}:

\begin{minted}{haskell}
codegenExp :: Exp -> CgM IR.Expr
codegenExp (EInt n)      = return (IR.CONST n)
codegenExp (EBool True)  = return (IR.CONST 1)
codegenExp (EBool False) = return (IR.CONST 0)
codegenExp (EString _)   = return (IR.CONST 0)
codegenExp (EVar v)      = return (IR.TEMP v)

codegenExp (e1 :+: e2)  = IR.BINOP IR.BAdd <$> codegenExp e1 <*> codegenExp e2
codegenExp (e1 :*: e2)  = IR.BINOP IR.BMul <$> codegenExp e1 <*> codegenExp e2
codegenExp (e1 :-: e2)  = IR.BINOP IR.BSub <$> codegenExp e1 <*> codegenExp e2
codegenExp (e1 :/: e2)  = IR.BINOP IR.BDiv <$> codegenExp e1 <*> codegenExp e2
codegenExp (e1 :<: e2)  = IR.BINOP IR.BLt  <$> codegenExp e1 <*> codegenExp e2
codegenExp (e1 :=: e2)  = IR.BINOP IR.BEq  <$> codegenExp e1 <*> codegenExp e2

codegenExp (e1 :&&: e2) = IR.BINOP IR.BAnd <$> codegenExp e1 <*> codegenExp e2
codegenExp (e1 :||: e2) = IR.BINOP IR.BOr  <$> codegenExp e1 <*> codegenExp e2

codegenExp (Not e) = IR.BINOP IR.BEq (IR.CONST 0) <$> codegenExp e
\end{minted}

Dois pontos merecem atenção:

\begin{itemize}
  \item \textbf{Booleanos.}  \texttt{true} compila para \texttt{CONST(1)} e
    \texttt{false} para \texttt{CONST(0)}.  Como os únicos valores booleanos
    são $0$ e $1$, os operadores bit a bit \texttt{AND} e \texttt{OR}
    coincidem com os conectivos lógicos \texttt{\&\&} e \texttt{||}.
  \item \textbf{Negação lógica.}  A expressão \texttt{not~e} é traduzida para
    $\mathbf{BINOP}(\mathtt{EQ},\; \mathbf{CONST}(0),\;
    \mathcal{E}\llbracket e \rrbracket)$:
    se $e$ avalia para $0$ (falso), o resultado é $1$; se avalia para $1$
    (verdadeiro), o resultado é $0$.
  \item \textbf{Strings.}  A IRT é inteiramente inteira; strings não têm
    representação direta e são compiladas para \texttt{CONST(0)} como
    aproximação conservadora.
\end{itemize}

\subsection{Exemplo: Fatorial Iterativo}\label{sec:ir-impl-example}

O programa TWhile abaixo computa $5!$ de forma iterativa:

\begin{minted}{text}
let n : int := 5;
let acc : int := 1;
while n > 0 {
    acc := acc * n;
    n := n - 1;
}
print acc;
\end{minted}

Após a invocação de \texttt{compileTWhile}, o código IRT gerado é (com rótulos
renomeados para clareza):

\begin{minted}{text}
func main() {
  MOVE(TEMP(n),   CONST(5))
  MOVE(TEMP(acc), CONST(1))
  LABEL(while_head)
  CJUMP(BINOP(LT, CONST(0), TEMP(n)), while_body, while_exit)
  LABEL(while_body)
    MOVE(TEMP(acc), BINOP(MUL, TEMP(acc), TEMP(n)))
    MOVE(TEMP(n),   BINOP(SUB, TEMP(n),   CONST(1)))
    JUMP(NAME(while_head))
  LABEL(while_exit)
  EXP(CALL(NAME(print), TEMP(acc)))
}
\end{minted}

A saída é uma função \texttt{main} sem parâmetros cuja sequência de
instruções espelha diretamente a estrutura do programa fonte.

\subsection{Integração no Pipeline}\label{sec:ir-impl-pipeline}

O executável \texttt{twhile} expõe a opção \texttt{--ir}:

\begin{minted}{text}
twhile --ir -f programa.tw
\end{minted}

O pipeline executa as seguintes etapas em ordem:
\begin{enumerate}
  \item \textbf{Análise sintática:} o parser produz a AST de TWhile.
  \item \textbf{Verificação de tipos:} o analisador semântico valida a AST
    e devolve o contexto de tipos final.
  \item \textbf{Geração de código:} \texttt{compileTWhile} converte a AST
    verificada em um \texttt{IR.Program}.
  \item \textbf{Impressão:} o módulo \texttt{IR.Frontend.Pretty.IRPretty}
    serializa o programa IRT para texto e o resultado é escrito num arquivo
    \texttt{.ir} com o mesmo nome base do arquivo fonte.
\end{enumerate}

O arquivo \texttt{.ir} gerado pode então ser interpretado diretamente pelo
executável \texttt{ir --interp}, fechando o ciclo
\emph{fonte} $\to$ \emph{IRT} $\to$ \emph{execução}.

% ============================================================
\section{Geração de Código para TImp: Funções e Registros}
\label{sec:ir-timp}
% ============================================================

A linguagem TImp estende TWhile com \textbf{declarações de funções} e
\textbf{tipos registro}. O \emph{backend} correspondente encontra-se em
\texttt{TImp.Backend.IR.TImpCodegen} e herda toda a estrutura de TWhile,
acrescentando:

\begin{enumerate}
  \item regras de tradução para acesso e atribuição de campos;
  \item uma regra de alocação para a construção de registros;
  \item a tradução de cada declaração de função para um fragmento
    \texttt{FuncDef} distinto;
  \item propagação do comando \texttt{return} via \texttt{RETURN}.
\end{enumerate}

% ------------------------------------------------------------
\subsection{Representação de Registros na Memória}
\label{sec:ir-record-layout}
% ------------------------------------------------------------

Registros são alocados no \textbf{heap} como blocos contíguos de palavras
de 64 bits (\(w = 8\) bytes). Os campos são numerados pela ordem de
declaração, a partir de 0. O campo de índice \(i\) de um registro cujo
endereço base está em \(\mathbf{TEMP}(x)\) está no endereço:

\[
  \mathit{addr}(x, i) = \mathbf{BINOP}(\mathtt{ADD},\; \mathbf{TEMP}(x),\;
  \mathbf{CONST}(i \cdot w))
\]

Para um registro \texttt{Point \{ x : int; y : int \}}, o campo \texttt{x}
tem índice 0 (deslocamento 0) e o campo \texttt{y} tem índice 1
(deslocamento 8). Não há cabeçalho de tipo: a distinção de tipo é
responsabilidade do sistema de tipos estático.

A alocação é realizada pela função embutida \texttt{alloc}, chamada com
o número de campos como argumento. O interpretador IRT implementa
\texttt{alloc} via um alocador \emph{bump-pointer} iniciado num endereço
fixo alto:

\begin{minted}{haskell}
callFunc _ (NAME "alloc") [n] = do
  ptr <- gets heapPtr
  modify (\s -> s { heapPtr = ptr + n * 8 })
  return ptr
\end{minted}

% ------------------------------------------------------------
\subsection{Tradução Dirigida por Sintaxe para Registros}
\label{sec:ir-record-sdt}
% ------------------------------------------------------------

\subsubsection*{Construção de Registro}

Seja \(R\) declarado com campos \((f_0,\ldots,f_{n-1})\) em ordem. A
expressão \(\mathbf{new}\; R\;\{f_0\!=\!e_0,\ldots,f_{n-1}\!=\!e_{n-1}\}\)
é traduzida para:

\[
  \mathcal{E}\llbracket \mathbf{new}\; R\;\{\overline{f_i\!=\!e_i}\} \rrbracket
  = \mathbf{ESEQ}(s_{\mathit{alloc}},\; \mathbf{TEMP}(t))
\]

onde \(t\) é um novo temporário e \(s_{\mathit{alloc}}\) é o bloco de
alocação e inicialização dos campos:

\[
  \begin{aligned}
    s_{\mathit{alloc}}
    &= \mathbf{SEQ}( \\
      &\quad \mathbf{MOVE}(\mathbf{TEMP}(t),\;
      \mathbf{CALL}(\mathbf{NAME}(\texttt{alloc}),\; \mathbf{CONST}(n))), \\
      &\quad \mathbf{SEQ}( \\
        &\qquad \mathbf{MOVE}(\mathbf{MEM}(\mathit{addr}(t,0)),\;
        \mathcal{E}\llbracket e_0 \rrbracket), \\
        &\qquad \cdots \\
        &\qquad \mathbf{MOVE}(\mathbf{MEM}(\mathit{addr}(t,n-1)),\;
    \mathcal{E}\llbracket e_{n-1} \rrbracket)))
  \end{aligned}
\]

O bloco \(s_{\mathit{alloc}}\) aloca
\(n\) palavras, armazena os valores dos campos em ordem e o
\(\mathbf{ESEQ}\) devolve \(\mathbf{TEMP}(t)\) como endereço do registro
recém-criado.

\subsubsection*{Acesso a Campo}

Dado que \(e\) tem tipo \(R\) e \(f\) tem índice \(i\) em \(R\):

\[
  \mathcal{E}\llbracket e.f \rrbracket =
  \mathbf{MEM}\!\left(
    \mathbf{BINOP}(\mathtt{ADD},\;
      \mathcal{E}\llbracket e \rrbracket,\;
    \mathbf{CONST}(i \cdot w))
  \right)
\]

Tal como o acesso a componente de tupla (Seção~\ref{sec:sdt-expr}), o
índice é convertido em deslocamento de byte.

\subsubsection*{Atribuição a Campo}

\[
  \mathcal{S}\llbracket x.f \mathbin{:=} e \rrbracket =
  \mathbf{MOVE}\!\left(
    \mathbf{MEM}\!\left(
      \mathbf{BINOP}(\mathtt{ADD},\;
        \mathbf{TEMP}(x),\;
      \mathbf{CONST}(i \cdot w))
    \right),\;
    \mathcal{E}\llbracket e \rrbracket
  \right)
\]

A semântica de valor de TImp é preservada: esta instrução \emph{não} cria
uma cópia do registro; atualiza o campo \(f\) no bloco já alocado, pois a
variável \(x\) detém o \emph{endereço} do registro. (A cópia ocorre apenas
  quando uma variável do tipo registro é passada como argumento ou atribuída
  a outra variável, ponto em que o \emph{front-end} pode inserir chamadas a
  \texttt{alloc}; na IRT, toda passagem de registro é por referência implícita
via o temporário que guarda o endereço.)

% ------------------------------------------------------------
\subsection{Tradução de Funções}
\label{sec:ir-timp-fn}
% ------------------------------------------------------------

Cada declaração de função
\(\mathbf{fn}\; f(\overline{x_i : \tau_i}) : \rho\; \{B\}\)
torna-se um \texttt{FuncDef} independente:

\[
  \mathcal{F}\llbracket \mathbf{fn}\; f(\overline{x_i:\tau_i}):\rho\;\{B\}
  \rrbracket
  = \mathtt{FuncDef}\; f\; [x_1,\ldots,x_n]\;
  (\mathit{foldStmts}(\mathcal{S}\llbracket B \rrbracket))
\]

Os nomes dos parâmetros formais \(x_i\) são usados diretamente como
temporários no corpo, o que é correto porque o interpretador IRT inicializa
\(\mathbf{TEMP}(x_i) = v_i\) no estado local de cada chamada:
não há prólogo de cópia de argumentos.

As regras de tradução para \textbf{return} são:

\[
  \mathcal{S}\llbracket \mathbf{return}\; e \mathtt{;} \rrbracket =
  \mathbf{RETURN}(\mathcal{E}\llbracket e \rrbracket)
  \qquad
  \mathcal{S}\llbracket \mathbf{return} \mathtt{;} \rrbracket =
  \mathbf{RETURN}()
\]

Chamadas de função como expressão traduzem-se como antes:
\[
  \mathcal{E}\llbracket f(e_1,\ldots,e_n) \rrbracket =
  \mathbf{CALL}(\mathbf{NAME}(f),\;
    \mathcal{E}\llbracket e_1 \rrbracket,\;\ldots,\;
  \mathcal{E}\llbracket e_n \rrbracket)
\]

% ------------------------------------------------------------
\subsection{Implementação em Haskell}
\label{sec:ir-timp-impl}
% ------------------------------------------------------------

\subsubsection*{Estado Estendido da Geração de Código}

O estado de TWhile é estendido com dois campos adicionais:

\begin{minted}{haskell}
data CgState = CgState
  { nextLabel  :: Int
  , nextTemp   :: Int
  , stmtAcc    :: [IR.Stmt]
  , recFields  :: Map Name [Field]   -- ordem de campos por tipo R
  , varTypes   :: Map Var Ty         -- tipo de cada variavel local
  }
\end{minted}

\begin{itemize}
  \item \texttt{recFields} mapeia o nome de cada tipo registro à lista
    ordenada de seus campos. Permite calcular o índice \(i\) de um campo
    \(f\) em \(O(\log n)\) com \texttt{elemIndex}.
  \item \texttt{varTypes} rastreia o tipo TImp de cada variável declarada.
    É necessário para traduzir \texttt{EField} e \texttt{SFieldAssign}: o
    nome do tipo registro deve ser recuperado para consultar
    \texttt{recFields}.
  \item \texttt{nextTemp} fornece novos temporários para a construção de
    registros, evitando colisões com os nomes de variáveis do programa.
\end{itemize}

O estado inicial é construído com \texttt{recFields} já preenchido e
\texttt{varTypes} vazio (as variáveis são inseridas à medida que as
declarações são processadas):

\begin{minted}{haskell}
initState :: Map Name [Field] -> CgState
initState rf = CgState 0 0 [] rf Map.empty
\end{minted}

\subsubsection*{Ponto de Entrada}

\begin{minted}{haskell}
compileTImp :: TImp -> Either String IR.Program
compileTImp (TImp decls stmts) = do
  let rf = Map.fromList
            [ (n, [f | FieldDecl f _ <- fields])
            | DRecord (RecordDecl n fields) <- decls ]
  funcDefs <- mapM (compileFuncDecl rf) [fd | DFunc fd <- decls]
  mainStmts <- runCgM rf (mapM_ codegenStmt stmts)
  let mainDef = IR.FuncDef "main" [] (foldStmts mainStmts)
  pure (funcDefs ++ [mainDef])
\end{minted}

O mapa \texttt{rf} é construído uma vez a partir das declarações de tipo.
Cada função é compilada independentemente por \texttt{compileFuncDecl};
os comandos de nível superior são colocados numa função \texttt{main}
implícita.

\subsubsection*{Compilação de Declarações de Função}

\begin{minted}{haskell}
compileFuncDecl :: Map Name [Field] -> FuncDecl
                -> Either String IR.FuncDef
compileFuncDecl rf (FuncDecl name params _ body) = do
  let paramNames = [v | Param v _ <- params]
      paramTypes = [(v, t) | Param v t <- params]
  stmts <- runCgM rf $ do
    mapM_ (uncurry trackVar) paramTypes  -- preenche varTypes
    mapM_ codegenStmt body
  pure (IR.FuncDef name paramNames (foldStmts stmts))
\end{minted}

Os tipos dos parâmetros são pré-carregados em \texttt{varTypes} antes de
iniciar a geração do corpo. Sem isso, um acesso \texttt{p.x} dentro da
função não conseguiria determinar o tipo de \texttt{p}.

\subsubsection*{Geração de Comandos: Novos Casos}

\textbf{Declaração de variável tipada} (rastreia o tipo para uso futuro):

\begin{minted}{haskell}
codegenStmt (SDecl v t e) = do
  ir <- codegenExp e
  emit (IR.MOVE (IR.TEMP v) ir)
  trackVar v t           -- insere v -> t em varTypes
\end{minted}

\textbf{Atribuição a campo} (implementa a regra $\mathcal{S}[x.f:=e]$):

\begin{minted}{haskell}
codegenStmt (SFieldAssign v f e) = do
  vts <- gets varTypes
  case Map.lookup v vts of
    Nothing -> throwError ("Unknown variable type for " ++ v)
    Just (TRecord rname) -> do
      i   <- fieldIndex rname f         -- indice i do campo f em R
      rhs <- codegenExp e
      let addr = IR.BINOP IR.BAdd (IR.TEMP v) (IR.CONST (i * wordSize))
      emit (IR.MOVE (IR.MEM addr) rhs)
    Just t -> throwError (v ++ " is not a record (type: " ++ show t ++ ")")
\end{minted}

\textbf{Retorno com e sem valor}:

\begin{minted}{haskell}
codegenStmt (SReturn Nothing)  = emit (IR.RETURN [])
codegenStmt (SReturn (Just e)) = do
  ir <- codegenExp e
  emit (IR.RETURN [ir])
\end{minted}

\textbf{Chamada de função como comando}:

\begin{minted}{haskell}
codegenStmt (SCall fname args) = do
  argIRs <- mapM codegenExp args
  emit (IR.EXP (IR.CALL (IR.NAME fname) argIRs))
\end{minted}

O valor de retorno é descartado via \texttt{EXP}, tal como em TWhile.

\subsubsection*{Geração de Expressões: Novos Casos}

\textbf{Acesso a campo} (implementa a regra $\mathcal{E}[e.f]$, restrito
a \texttt{EVar} para simplificar):

\begin{minted}{haskell}
codegenExp (EField (EVar v) f) = do
  vts <- gets varTypes
  case Map.lookup v vts of
    Nothing -> throwError ("Unknown variable type for " ++ v)
    Just (TRecord rname) -> do
      i <- fieldIndex rname f
      let addr = IR.BINOP IR.BAdd (IR.TEMP v)
                   (IR.CONST (i * wordSize))
      pure (IR.MEM addr)
    Just t -> throwError (v ++ " has non-record type " ++ show t)

codegenExp (EField _ _) =
  throwError "Field access on non-variable not supported in IR codegen"
\end{minted}

A restrição a \texttt{EVar} é conservadora mas correta para todos os
programas produzidos pelo \emph{front-end}: o verificador de tipos garante
que o receptor de um acesso a campo é bem tipado, e a expressão mais
comum é de facto uma variável. Expressões aninhadas como
\texttt{f().x} exigiriam um temporário intermediário.

\textbf{Construção de registro} (implementa a regra
$\mathcal{E}[\mathbf{new}\;R\{\ldots\}]$):

\begin{minted}{haskell}
codegenExp (ENew rname initFields) = do
  rf <- gets recFields
  case Map.lookup rname rf of
    Nothing   -> throwError ("Unknown record type: " ++ rname)
    Just flds -> do
      let n = length flds
      t <- freshTemp
      -- avalia campos na ordem de declaracao (nao na ordem fornecida)
      fieldVals <- mapM (lookupFieldVal initFields) flds
      let allocStmt  = IR.MOVE (IR.TEMP t)
                         (IR.CALL (IR.NAME "alloc") [IR.CONST n])
          fieldStmts = [ IR.MOVE
                           (IR.MEM (IR.BINOP IR.BAdd
                             (IR.TEMP t) (IR.CONST (i * wordSize)))) v
                       | (i, v) <- zip [0..] fieldVals ]
      pure (IR.ESEQ (foldStmts (allocStmt : fieldStmts)) (IR.TEMP t))
  where
    lookupFieldVal pairs f =
      case lookup f pairs of
        Nothing -> throwError ("Missing field " ++ f)
        Just e  -> codegenExp e
\end{minted}

Três aspectos merecem atenção:

\begin{itemize}
  \item Os campos são avaliados na \textbf{ordem de declaração} de \(R\)
    (lista \texttt{flds}), não na ordem em que o programador os escreveu.
    Isso garante que o deslocamento \(i \cdot w\) está correto
    independentemente da ordem de inicialização.
  \item O temporário \(t\) é gerado por \texttt{freshTemp} (contador
    \texttt{nextTemp}), que produz nomes como \texttt{\_t0}, \texttt{\_t1},
    distintos dos nomes de variáveis do programa.
  \item A chamada a \texttt{alloc} recebe \texttt{CONST(n)}, o número
    de campos, e não o número de bytes. O interpretador IRT multiplica
    por \(w\) internamente.
\end{itemize}

% ------------------------------------------------------------
\subsection{Exemplo Completo}
\label{sec:ir-timp-example}
% ------------------------------------------------------------

Considere o programa TImp:

\begin{minted}{text}
record Point { x : int; y : int; }
fn soma(p : Point) : int {
  return p.x + p.y;
}
var q : Point = new Point { x = 3, y = 4 };
print soma(q);
\end{minted}

Após \texttt{compileTImp}, o código IRT gerado é:

\begin{minted}{text}
func soma(p) {
  RETURN(BINOP(ADD,
    MEM(BINOP(ADD, TEMP(p), CONST(0))),   -- p.x  (indice 0)
    MEM(BINOP(ADD, TEMP(p), CONST(8))))) -- p.y  (indice 1)
}
func main() {
  SEQ(
    MOVE(TEMP(q),                         -- var q = new Point{x=3,y=4}
      ESEQ(
        SEQ(MOVE(TEMP(_t0), CALL(NAME(alloc), CONST(2))),
            SEQ(MOVE(MEM(BINOP(ADD, TEMP(_t0), CONST(0))), CONST(3)),
                MOVE(MEM(BINOP(ADD, TEMP(_t0), CONST(8))), CONST(4)))),
        TEMP(_t0))),
    EXP(CALL(NAME(print),                 -- print soma(q)
          CALL(NAME(soma), TEMP(q)))))
}
\end{minted}

Pontos a observar:

\begin{itemize}
  \item A função \texttt{soma} produz um \texttt{FuncDef} independente com
    parâmetro \texttt{p}. Os acessos \texttt{p.x} e \texttt{p.y} tornam-se
    \texttt{MEM} com deslocamentos 0 e 8.
  \item O registro é alocado com \texttt{alloc(2)} e o endereço base fica
    em \texttt{TEMP(\_t0)}. Os dois campos são inicializados em seguida.
  \item A chamada \texttt{soma(q)} passa o endereço do registro (o valor
    inteiro em \texttt{TEMP(q)}) como argumento; dentro de \texttt{soma},
    \texttt{TEMP(p)} contém esse endereço.
  \item O resultado de \texttt{soma(q)} é usado diretamente como argumento
    de \texttt{print}, sem temporário intermediário (a canonicalização
    extrairia essa chamada embutida se necessário).
\end{itemize}

% ============================================================
\section{Conclusão}\label{sec:ir-conclusao}

A representação intermediária em árvore (IRT) introduzida neste capítulo
satisfaz os critérios de simplicidade, independência de máquina e suporte a
transformações enunciados na Seção~\ref{sec:ir-desiderata}.
A semântica big-step formaliza o comportamento de cada construtor,
permitindo raciocinar com rigor sobre a correção das traduções.
As regras de tradução dirigida por sintaxe cobrem todas as construções de
TWhile e TImp e são estruturalmente indutivas sobre a AST, o que facilita
tanto a implementação quanto a prova de corretude.

Os próximos passos na cadeia de compilação são:
\begin{itemize}
  \item \textbf{Canonicalização e formação de blocos básicos:}
    preparar a IRT linearizada para análise de fluxo.
  \item \textbf{Otimizações sobre a RI:}
    propagação de constantes, eliminação de código morto,
    movimentação de código invariante.
  \item \textbf{Seleção de instruções:}
    mapear construções da IRT para instruções da arquitetura-alvo
    (e.g., x86-64 ou WebAssembly).
  \item \textbf{Alocação de registradores:}
    decidir, para cada temporário, se ele vive em um registrador físico
    ou em uma posição da pilha.
\end{itemize}
